\ifdefined\gkver\else\def\gkver{student}\fi
\documentclass[\gkver]{gaokaozhenti}
\newcommand{\ccnum}[1]{{\fontspec{Noto Sans CJK JP}#1}}
\begin{document}

\chapter{2015年广东理综}
%% paperType: 理综物理部分

\begin{enumerate}
\item
%% number: 13
%% paperName: 2015年广东理综第 13 题 4 分
%% typeId: 1
%% type: 单选题
%% chapter: 直线运动
%% point: 运动图象
%% method: 无
%% score: 4
%% degree: 650
%% duplicateId: 0
%% body:
甲、乙两人同时同地出发骑自行车做直线运动，前 1 小时内的位移-时间图像如图所示。下列表述正确的是\xzanswer{B}

\onepicture[h=3.4cm]{figs/13.png}

\fourchoices[answer=B]
{0.2\textasciitilde{}0.5 小时内，甲的加速度比乙的大}
{0.2\textasciitilde{}0.5 小时内，甲的速度比乙的大}
{0.6-0.8 小时内，甲的位移比乙的小}
{0.8 小时内，甲、乙骑行的路程相等}

%% answer: B
%% memo:
\memoanswer{
	由位移-时间图知 0.2\textasciitilde{}0.5 小时内，甲、乙都做匀速直线运动，加速度都为 0，故 A 错误；图线斜率表示速度大小，所以甲的速度大于乙的速度，B 正确；0.6\textasciitilde{}0.8 小时内，甲的位移是 $-5\Um$，乙的位移是 $-3\Um$，C 错误；0.8 小时内甲乙的位移相同，但是路程不一样，D 错误，选项 B 正确。
}

\item
%% number: 14
%% paperName: 2015年广东理综第 14 题 4 分
%% typeId: 1
%% type: 单选题
%% chapter: 直线运动
%% point: 运动的合成
%% method: 无
%% score: 4
%% degree: 450
%% duplicateId: 0
%% body:
如图所示，帆板在海面上以速度 $v$ 朝正西方向运动，帆船以速度 $v$ 朝正北方向航行，以帆板为参照物\xzanswer{D}

\onepicture[h=3.2cm]{figs/14.png}

\fourchoices[answer=D]
{帆船朝正东方向航行，速度大小为 $v$}
{帆船朝正西方向航行，速度大小为 $v$}
{帆船朝南偏东 $45^{\circ}$ 方向航行，速度大小为 $\sqrt{2}v$}
{帆船朝北偏东 $45^{\circ}$ 方向航行，速度大小为 $\sqrt{2}v$}

%% answer: D
%% memo:
\memoanswer{
	以帆板为参照物，就把帆板看成是静止的，则帆船在水平方向以速度 $v$ 向东运动，竖直方向以速度 $v$ 向北运动，根据矢量合成的平行四边形定则可以求得帆船以帆板为参照物是以 $v$ 的速度向北偏东 $45^{\circ}$ 运动，故选项 D 正确。
}

\item
%% number: 15
%% paperName: 2015年广东理综第 15 题 4 分
%% typeId: 1
%% type: 单选题
%% chapter: 交流电
%% point: 变压器
%% method: 无
%% score: 4
%% degree: 650
%% duplicateId: 0
%% body:
如图为加热装置的示意图，使用电阻丝加热导气管，视变压器为理想变压器，原线圈接入电压有效值恒定的交流电并保持匝数不变，调节触头 P，使输出电压有效值由 $220\UV$ 降至 $110\UV$。调节前后\xzanswer{C}

\onepicture[h=3.2cm]{figs/15.png}

\fourchoices[answer=C]
{副线圈中的电流比为 1:2}
{副线圈输出功率比为 2:1}
{副线圈的接入匝数比为 2:1}
{原线圈输入功率比为 1:2}

%% answer: C
%% memo:
\memoanswer{
	原线圈的输入电压和匝数不变，根据输出电压的有效值由 $220\UV$ 降到 $110\UV$，可以知道副线圈的匝数变为原来的 $\frac{1}{2}$，C 正确；根据 $P=\frac{U^{2}}{R}$ 可以知道，副线圈的输出功率变为原来的 $\frac{1}{4}$，同样原线圈的输入功率也变为原来的 $\frac{1}{4}$，B、D 错误；又 $I=\frac{U}{R}$，副线圈的电流应该变为原来的 $\frac{1}{2}$，A 错误，选项 C 正确。
}

\item
%% number: 16
%% paperName: 2015年广东理综第 16 题 4 分
%% typeId: 1
%% type: 单选题
%% chapter: 磁场
%% point: 带电粒子在磁场中的运动
%% method: 无
%% score: 4
%% degree: 650
%% duplicateId: 0
%% body:
在同一匀强磁场中，$\alpha$ 粒子（\ce{^{4}_{2}He}）和质子（\ce{^{1}_{1}H}）做匀速圆周运动，若他们的动量大小相等，则 $\alpha$ 粒子和质子\xzanswer{B}

\fourchoices[answer=B]
{运动半径之比是 2:1}
{运动周期之比是 2:1}
{运动速度大小之比是 4:1}
{受到的洛伦兹力之比是 2:1}

%% answer: B
%% memo:
\memoanswer{
	质量之比是 4:1，电荷量之比是 2:1，根据动量相等可以知道其速度之比是 1:4，洛伦兹力为 $qvB$，所以洛伦兹力之比为 1:2，C、D 错误；由 $qvB=m\frac{v^{2}}{R}$ 及 $T=\frac{2\pi R}{v}$ 可得：半径之比是 1:2，周期之比是 2:1，A 错误、B 正确，选项 B 正确。
}

\item
%% number: 17
%% paperName: 2015年广东理综第 17 题 6 分
%% typeId: 2
%% type: 多选题
%% chapter: 气体性质
%% point: 热力学第一定律
%% method: 无
%% score: 6
%% degree: 650
%% duplicateId: 0
%% body:
如图为某实验器材的结构示意图，金属内筒和隔热外筒间封闭了一定体积的空气，内筒中有水。在水加热升温的过程中，被封闭的空气\xzanswer{AB}

\onepicture[h=3.4cm]{figs/17.png}

\fourchoices[answer=AB]
{内能增大}
{压强增大}
{分子间引力和斥力都减小}
{所有分子运动速率都增大}

%% answer: AB
%% memo:
\memoanswer{
	空气体积不变，故做功为 0，即 $W=0$。由于外壁隔热，无热传递。水加热升温会吸收热量，故 $Q>0$，由热力学第一定律 $\Delta U=Q+W$ 得：$\Delta U>0$，内能增加，A 正确；温度升高，分子热运动变剧烈，对器壁撞击加强，故压强变大，B 正确；由于体积、质量不变，分子间距离不变，故作用力大小不变，C 错误；温度升高，分子的平均热运动速率加快，不代表每一个都加快，D 错误，选项 AB 正确。
}

\item
%% number: 18
%% paperName: 2015年广东理综第 18 题 6 分
%% typeId: 2
%% type: 多选题
%% chapter: 原子和原子核
%% point: 核聚变
%% method: 无
%% score: 6
%% degree: 650
%% duplicateId: 0
%% body:
科学家使用核反应获取氚，再利用氘和氚的核反应获得能量，核反应方程分别为：$\ce{X + Y -> ^{4}_{2}He + ^{3}_{1}H} + 4.9\UMeV$ 和 $\ce{^{2}_{1}H + ^{3}_{1}H -> ^{4}_{2}He + X} + 17.6\UMeV$，下列表述正确的有\xzanswer{AD}

\fourchoices[answer=AD]
{X 是中子}
{Y 的质子数是 3，中子数是 6}
{两个核反应都没有质量亏损}
{氘和氚的核反应是核聚变反应}

%% answer: AD
%% memo:
\memoanswer{
	由核反应过程质量守恒和电荷量守恒可知 X 是中子，Y 的质子数为 6，中子数为 6，A 正确、B 错误；两个核反应方程为核聚变反应，聚变放出能量，一定有质量亏损，C 错误、D 正确，选项 AD 正确。
}

\item
%% number: 19
%% paperName: 2015年广东理综第 19 题 6 分
%% typeId: 2
%% type: 多选题
%% chapter: 力物体的平衡
%% point: 多力平衡
%% method: 无
%% score: 6
%% degree: 450
%% duplicateId: 0
%% body:
如图所示，三条绳子的一端都系在细直杆顶端。另一端都固定在水平地面上。将杆竖直紧压在地面上。若三条绳长度不同，下列说法正确的有\xzanswer{BC}

\onepicture[h=3.4cm]{figs/19.png}

\fourchoices[answer=BC]
{三条绳中的张力都相等}
{杆对地面的压力大于自身重力}
{绳子对杆的拉力在水平方向的合力为零}
{绳子拉力的合力与杆的重力是一对平衡力}

%% answer: BC
%% memo:
\memoanswer{
	由受力分析可知：三个绳子长度不同，即与地面夹角不等，故其张力大小不等，A 错误；由于支持力与重力和绳子向下的分力平衡，故对地面的压力应与重力和绳子向下分力的和等大，B 正确；杆在水平方向处于平衡状态，故三绳水平方向合力为零，C 正确；绳的拉力，重力和杆受到的支持力三力平衡，D 错误，选项 BC 正确。
}

\item
%% number: 20
%% paperName: 2015年广东理综第 20 题 6 分
%% typeId: 2
%% type: 多选题
%% chapter: 曲线运动
%% point: 人造地球卫星
%% method: 无
%% score: 6
%% degree: 450
%% duplicateId: 0
%% body:
在星球表面发射探测器，当发射速度为 $v$ 时，探测器可绕星球表面做匀速圆周运动；当发射速度为 $2v$ 时，可摆脱星球引力束缚脱离该星球。已知地球、火星两星球的质量比约为 10:1、半径比约为 2:1，下列说法正确的有\xzanswer{BD}

\fourchoices[answer=BD]
{探测器的质量越大，脱离星球所需要的发射速度越大}
{探测器在地球表面受到的引力比在火星表面的大}
{探测器分别脱离两星球所需要的发射速度相等}
{探测器脱离星球的过程中，势能逐渐增大}

%% answer: BD
%% memo:
\memoanswer{
	根据 $G\frac{Mm}{R^{2}}=m\frac{v^{2}}{R}$ 得发射速度 $v=\sqrt{\frac{GM}{R}}$，脱离星球所需要的发射速度为 $v'=2v=2\sqrt{\frac{GM}{R}}$，与探测器的质量无关，A 错误；由 $F=G\frac{Mm}{R^{2}}$，则 $\frac{F_{\text{地}}}{F_{\text{星}}}=\frac{M_{\text{地}}R_{\text{星}}^{2}}{M_{\text{星}}R_{\text{地}}^{2}}=\frac{5}{2}$，地球表面的引力大于火星表面的引力，B 正确；由脱离星球发射速度 $v'=2\sqrt{\frac{GM}{R}}$ 可得 $\frac{v_{\text{地}}'}{v_{\text{星}}'}=\sqrt{\frac{M_{\text{地}}R_{\text{星}}}{M_{\text{星}}R_{\text{地}}}}=\sqrt{5}$，C 错误；脱离星球的过程中，高度在增加，势能增大，D 正确。选项 BD 正确。
}

\item
%% number: 21
%% paperName: 2015年广东理综第 21 题 6 分
%% typeId: 2
%% type: 多选题
%% chapter: 电场
%% point: 带电粒子的轨迹类问题
%% method: 无
%% score: 6
%% degree: 450
%% duplicateId: 0
%% body:
如图所示的水平匀强电场中，将两个带电小球 M 和 N 分别沿图示路径移动到同一水平线上的不同位置，释放后，M、N 保持静止，不计重力，则\xzanswer{BD}

\onepicture[h=2.8cm]{figs/21.png}

\fourchoices[answer=BD]
{M 的带电量比 N 的大}
{M 带负电荷、N 带正电荷}
{静止时 M 受到的合力比 N 的大}
{移动过程中匀强电场对 M 做负功}

%% answer: BD
%% memo:
\memoanswer{
	释放后，M、N 保持静止，两小球各自受电场力和两小球间的库仑力等大反向，故两小球所受电场力大小相等，在匀强电场中：$F=qE$，则 $q$ 相等，A 错误；

	若 M 带负电荷、N 带正电荷，则 M 受到向左的电场力，向右的库仑力可平衡，N 受到向左的库仑力和向右电场力可平衡，B 正确；

	静止时，受合力为零，C 错误；M 带负电荷受到向左的电场力，向右运动电场力做负功，D 正确。

	选项 D 正确。
}

\item
%% number: 34
%% paperName: 2015年广东理综第 34 题 10 分
%% typeId: 4
%% type: 实验题
%% chapter: 电路
%% point: 伏安法测电阻
%% method: 无
%% score: 10
%% degree: 450
%% duplicateId: 0
%% body:
（1）（8 分）某同学使用打点计时器测量当地的重力加速度。

\begin{steps}
	\item
	请完成以下主要实验步骤：按图（a）安装实验器材并连接电源；竖直提起系有重物的纸带，使重物\tkanswer[1.2cm]{靠近}（填“靠近”或“远离”）计时器下端；\tkanswer[1.3cm]{接通电源}，\tkanswer[1.4cm]{释放纸带}，使重物自由下落；关闭电源，取出纸带；换新纸带重复实验。

	\item
	图（b）和（c）是实验获得的两条纸带，应选取\tkanswer[0.8cm]{b}（填“b”或“c”）来计算重力加速度。在实验操作和数据处理都正确的情况下，得到的结果仍小于当地重力加速度，主要原因是空气阻力和\tkanswer[1.6cm]{摩擦阻力}。
\end{steps}

\twopicture[showlabel=false, hA=4.0cm, hB=2.4cm]
{figs/34a.png}
{figs/34b.png}

（2）（10 分）某实验小组研究两个未知元件 X 和 Y 的伏安特性，使用的器材包括电压表（内阻约为 $3\UkO$），电流表（内阻约为 $1\UO$），定值电阻等。

\threepicture[showlabel=false, hA=2.6cm, hB=2.4cm, hC=2.4cm]
{figs/34c.png}
{figs/34d.png}
{figs/34e.png}

\begin{steps}
	\item
	使用多用电表粗测元件 X 的电阻。选择“$\times1$”欧姆挡测量，示数如图 1（a）所示，读数为\tkanswer[1.4cm]{10}$\UO$。据此应该选择图 1 中的\tkanswer[0.8cm]{b}（填“b”或“c”）的电路进行实验。

	\item
	连接所选电路，闭合 S；滑动变阻器的滑片 P 从左向右滑动，电流表的示数逐渐\tkanswer[1.2cm]{增大}（填“增大”或“减小”）；依次记录电流及相应的电压；将元件 X 换成 Y，重复实验。

	\item
	如图 2（a）是根据实验数据作出的 $U-I$ 图线，由图可判断元件\tkanswer[0.8cm]{Y}（填“X”或“Y”）是非线性元件。

	\item
	该小组还借助 X 和 Y 中的线性元件和阻值 $R=21\UO$ 的定值电阻，测量待测电池的电动势 $E$ 和内阻 $r$，电路如图 2（b）所示。闭合 $S_{1}$ 和 $S_{2}$，电压表读数为 $3.00\UV$；断开 $S_{2}$，读数为 $1.00\UV$。利用图（a）可算得 $E=$\tkanswer[1.2cm]{3.2}$\UV$，$r=$\tkanswer[1.2cm]{0.50}$\UO$（结果均保留两位有效数字，视电压表为理想电压表）。
\end{steps}

\twopicture[showlabel=false, hA=3.2cm, hB=2.8cm, align=right]
{figs/34f.png}
{figs/34g.png}

%% answer:
\jdanswer{
\begin{enumerate}
	\item
	（1）靠近，接通电源，释放纸带；（2）b，摩擦阻力。

	\item
	（1）$10\UO$，b；（2）增大；（3）Y；（4）$3.2\UV$，$0.50\UO$。
\end{enumerate}
}
%% memo:
\memoanswer{
	（1）①为了使实验数据能更好地记录在纸带上及减小实验误差，重物应靠近打点计时器；②由图可看出（b）纸带做加速运动、（c）纸带包含加速和减速过程，由自由落体运动是匀加速运动，故选（b）来计算加速度，由于空气阻力和摩擦阻力的作用使测得的加速度小于重力加速度，测量结果偏小。

	（2）①由题图读得：$10\times1\UO=10\UO$，为小电阻，为了减小误差应采用电流表外接法，即图 15（b）接法；②闭合 S，滑动变阻器的滑片 P 从左向右滑动，使测量电路电压增大，由欧姆定律得电流增大；③根据线性元件的 $U-I$ 图线为直线，非线性元件的 $U-I$ 图线为曲线可得：Y 是非线性元件；④由图 16（a）可得线性元件 X 的阻值为 $10\UO$，根据闭合电路欧姆定律及题意有：$E=3.00\UV+\frac{3.00}{10}r$ 和 $E=1.00\UV+\frac{1.00}{10}(21+r)$，解得 $E=3.15\UV\approx3.2\UV$，$r=0.50\UO$。
}

\item
%% number: 35
%% paperName: 2015年广东理综第 35 题 18 分
%% typeId: 6
%% type: 计算题
%% chapter: 电磁感应
%% point: 电磁感应中图象
%% method: 无
%% score: 18
%% degree: 450
%% duplicateId: 0
%% body:
如图（a）所示，平行长直金属导轨水平放置，间距 $L=0.4\Um$。导轨右端接有阻值 $R=1\UO$ 的电阻，导体棒垂直放置在导轨上，且接触良好。导体棒及导轨的电阻均不计，导轨间正方形区域 abcd 内有方向竖直向下的匀强磁场，bd 连线与导轨垂直，长度也为 $L$，从 0 时刻开始，磁感应强度 $B$ 的大小随时间 $t$ 变化，规律如图（b）所示；同一时刻，棒从导轨左端开始向右匀速运动，$1\Us$ 后刚好进入磁场。若使棒在导轨上始终以速度 $v=1\Ums$ 做直线运动，求：

\begin{enumerate}
	\item
	棒进入磁场前，回路中的电动势 $E$；

	\item
	棒在运动过程中受到的最大安培力 $F$，以及棒通过三角形 abd 区域时电流 $i$ 与时间 $t$ 的关系式。
\end{enumerate}

\twopicture[showlabel=false, hA=3.4cm, hB=3.0cm, align=right]
{figs/35a.png}
{figs/35b.png}

%% answer:
\jdanswer{
\begin{enumerate}
	\item
	$0.04\UV$

	\item
	$0.04\UN$，$i=\frac{2Bv^{2}}{R}t$
\end{enumerate}
}
%% memo:
\memoanswer{
	（1）棒在进入磁场前，棒没有做切割磁感线，但磁场的强弱发生变化，导致磁通量发生变化。

	abcd 的面积 $S=\frac{1}{2}L^{2}$\ccnum{①}，

	$E=n\frac{\Delta\Phi}{\Delta t}=n\frac{\Delta BS}{\Delta t}$\ccnum{②}，

	由\ccnum{①}\ccnum{②}联立得：$E=0.04\UV$。

	（2）棒进入磁场中后，做切割磁感线运动，当棒到达 bd 时，产生的感应电流最大，同时切割长度最大，到达 bd 时，产生的感应电动势 $E=BL_{bd}v$\ccnum{③}，

	产生的感应电流 $I=\frac{E}{R}$\ccnum{④}，

	所受最大安培力 $F=BIL_{bd}$\ccnum{⑤}，

	由\ccnum{③}\ccnum{④}\ccnum{⑤}联立得：$F=0.04\UN$。

	棒通过三角形 abd 区域时，切割的长度 $l=2vt$\ccnum{⑥}，

	产生的感应电动势 $E'=Blv$\ccnum{⑦}，

	感应电流 $i=\frac{E'}{R}$\ccnum{⑧}，

	由\ccnum{⑥}\ccnum{⑦}\ccnum{⑧}联立感应电流为：$i=\frac{2Bv^{2}}{R}t$。
}

\item
%% number: 36
%% paperName: 2015年广东理综第 36 题 18 分
%% typeId: 6
%% type: 计算题
%% chapter: 运动和力
%% point: 动量和能量
%% method: 无
%% score: 18
%% degree: 450
%% duplicateId: 0
%% body:
如图所示，一条带有圆轨道的长轨道水平固定，圆轨道竖直，底端分别与两侧的直轨道相切，半径 $R=0.5\Um$。物块 A 以 $v_{0}=6\Ums$ 的速度滑入圆轨道，滑过最高点 Q，再沿圆轨道滑出后，与直轨上 P 处静止的物块 B 碰撞，碰后粘在一起运动，P 点左侧轨道光滑，右侧轨道呈粗糙段，光滑段交替排列，每段长度都为 $L=0.1\Um$。物块与各粗糙段间的动摩擦因素都为 $\mu=0.1$，A、B 的质量均为 $m=1\Ukg$（重力加速度 $g$ 取 $10\Umsq$；A、B 视为质点，碰撞时间极短）。

\begin{enumerate}
	\item
	求 A 滑过 Q 点时的速度大小 $v$ 和受到的弹力大小 $F$；

	\item
	若碰后 AB 最终停止在第 $k$ 个粗糙段上，求 $k$ 的数值；

	\item
	求碰后 AB 滑至第 $n$ 个（$n<k$）光滑段上的速度 $v_{n}$ 与 $n$ 的关系式。
\end{enumerate}

\onepicture[align=right, h=3.6cm]{figs/36.png}

%% answer:
\jdanswer{
\begin{enumerate}
	\item
	$4\Ums$，$102\UN$

	\item
	45

	\item
	$v_{n}=\sqrt{v_{1}^{2}-2\mu gnL}$
\end{enumerate}
}
%% memo:
\memoanswer{
	（1）物块 A 滑入圆轨道到达 Q 的过程中机械能守恒，根据机械能守恒：

	$\frac{1}{2}mv_{0}^{2}=\frac{1}{2}mv^{2}+2mgR$\ccnum{①}，

	物块 A 在做圆周运动：$F_{N}+mg=m\frac{v^{2}}{R}$\ccnum{②}，

	由\ccnum{①}\ccnum{②}联立得：$v=4\Ums$，

	$F_{N}=102\UN$，方向向下\ccnum{③}。

	（2）在与 B 碰撞前，系统机械能守恒，所以与 B 碰前 A 的速度为 $6\Ums$，A 和 B 在碰撞过程中动量守恒：

	$m_{A}v_{0}=(m_{A}+m_{B})v_{1}$\ccnum{④}，

	AB 碰后向右滑动，由动能定理得：

	$-\mu(m_{A}+m_{B})gs=0-\frac{1}{2}(m_{A}+m_{B})v_{1}^{2}$\ccnum{⑤}，

	由\ccnum{④}\ccnum{⑤}联立得 $s=4.5\Um$，

	$\therefore k=\frac{s}{L}=45$。

	（3）碰后 AB 滑至第 $n$ 个光滑段上的速度 $v_{n}$，由动能定理得：

	$-\mu(m_{A}+m_{B})gnL=\frac{1}{2}(m_{A}+m_{B})v_{n}^{2}-\frac{1}{2}(m_{A}+m_{B})v_{1}^{2}$\ccnum{⑥}，

	解得：$v_{n}=\sqrt{v_{1}^{2}-2\mu gnL}$。
}

\end{enumerate}

\end{document}
